%%
% BIThesis 本科毕业设计论文模板 —— 使用 XeLaTeX 编译 The BIThesis Template for Undergraduate Thesis
% This file has no copyright assigned and is placed in the Public Domain.
%%

\begin{conclusion}
本文围绕"面向桌面操作的视觉—语言—动作机械臂抓取与控制"这一课题，在 Franka FR3 七自由度机械臂与二指夹爪平台上，完成了从实时控制接入、LeRobot 通用接口扩展、示教数据采集与规范化，到 $\pi_{0.5}$ 模型迁移微调、变体改进与实机闭环部署的全栈技术链路。现按设计任务逐项总结完成情况。

T1——弗兰卡实时控制与 LeRobot 接口扩展（第二章）：
设计并实现了四层系统架构（实时执行层—网络服务层—LeRobot 抽象层—遥操作与数据层）。在实时侧部署基于 libfranka 与 Ruckig 的实时控制服务进程，通过 TCP 套接字与上位机解耦，保证 $1\,\mathrm{kHz}$ 底层控制与 $\geqslant 30\,\mathrm{Hz}$ 指令刷新频率。在 LeRobot 框架中实现了机器人抽象及夹爪复合配置，使本平台可被社区录制工具直接消费。在遥操作层面实现了三维鼠标笛卡尔增量与同构主从关节空间两种示教入口，并设计了交互式两阶段配对标定流程——第一阶段以一次对准姿态自动解算七轴缩放比例与偏置，第二阶段测量夹爪物理量程——将标定结果写入配置文件实现参数版本化。经标定后，主从映射的关节空间偏差均值控制在 $1^\circ$ 以内，示教轨迹的可重复性显著提升。

T2——桌面抓取示教数据采集与规范化（第二章）：
以"抓取绿色苹果并放置至桌面左上角"为任务实例，采用双相机布局（第三人称固定相机与腕部手眼相机），通过同构主从臂完成精细示教，按 $10\,\mathrm{Hz}$ 帧率同步录制了桌面抓取数据集。数据集严格遵循 LeRobot v2.1 目录规范：包含 30 条有效轨迹、5,512 帧、60 个视频轨道；元数据目录中定义了任务语义、轨迹信息及 8 维状态/动作特征描述，数据目录以列式存储格式存放帧级表格，视频目录存放双路 RGB 视频。数据集字段编排与 OpenPI 训练接口直接对齐，可被 LeRobot 数据加载器直接消费。

T3——$\pi_{0.5}$ 模型迁移与变体改进（第三章、第四章）：
将双视角观测、8 维状态/动作与自然语言提示映射至 OpenPI 训练接口，在加载 $\pi_{0.5}$ 预训练权重的基础上完成全量微调基线（成功率约 $80\%$）。在此基础上，针对桌面任务小样本、接触敏感的特性，提出并实现了两项模型结构改进。第一项是冻结骨干轻量化变体：利用参数路径过滤机制精确冻结 VLM 骨干而仅微调动作专家，训练显存占用从约 160 GB 降至约 60 GB（批次大小为 4），成功率提升至 $90\%$。第二项是短期记忆变体：在视觉编码器中引入空间-时间可分离注意力与时间位置编码，当前帧保留完整空间细节、历史帧压缩为少量池化摘要 token，为模型提供短期视觉记忆；成功率达到 $100\%$，并表现出自主失败恢复能力——在初次夹取失败后能够自发重新搜索目标并再次尝试抓取。

T4——实机闭环部署与在线执行策略对比（第三章）：
针对同一模型权重在不同部署方式下可能产生截然不同实机表现的问题，设计并实施了受控对比实验。在同一训练完成的 $\pi_{0.5}$ 权重下，分别采用异步远端推理（执行完整动作块，偏长时开环）和分段同步闭环（每次仅执行动作块的一小段，约 $0.27\,\mathrm{s}$ 后重新观测再推理）进行 10 次桌面苹果抓取测试。实验结果表明：异步开环方式成功率仅约 $5\%$，失败集中在夹取接触瞬间的滑移与掉落；分段同步闭环方式成功率接近 $100\%$，显著增强了接触阶段的在线修正与失败恢复能力。该对比量化了在线闭环策略对接触敏感任务的决定性作用——成功率差异超过 $90$ 个百分点——验证了研究假设 H2，为 VLA 在类似桌面操作任务中的部署提供了可操作的工程准则。

创新点总结。本文的创新之处在于工程与方法的紧密结合。首先，弗兰卡--LeRobot 一体化接口与可版本化标定链路，将"通用硬件—标准数据—可复用训练栈"打通为可复现流程。其次，基于参数路径过滤的冻结骨干与基于空间-时间可分离注意力机制的短期记忆两项模型结构改进，通过配置驱动设计实现了训练显存效率与任务成功率的双重提升。再次，通过受控部署对比实验，首次在本平台桌面抓取场景中系统量化了闭环执行策略对 VLA 动作块预测模型的性能影响量级，揭示了"部署策略与模型结构同等重要"的工程规律。

不足与展望。受时间与数据规模限制，本文任务分布较窄（单一绿色苹果抓取），未系统评估跨物体、跨光照、跨台面布局的泛化能力。冻结骨干与短期记忆变体的超参数配置尚未充分网格化对比。部署对比仅在本平台上完成，尚未验证结论在其他机械臂平台上的可迁移性。后续工作可沿以下方向展开：扩展多物体、多姿态、多场景的示教数据集，结合数据增强提升策略泛化性；对模型变体进行更系统的超参数扫描与消融实验；将分段闭环部署策略推广至更复杂的多阶段长程任务；探索仿真到真实迁移与安全性约束下的在线策略切换机制。

总体而言，本文通过一条可复现的全栈技术路线证明了：先进 VLA 模型在静态桌面平台上的实用价值，不仅取决于模型本身的规模与预训练数据，更取决于与真实执行链路、标准化数据规范和在线闭环策略形成一致、可诊断的整体系统设计。这一结论对同类实验室平台引进与落地 VLA 技术具有参考意义。
\end{conclusion}
